\documentclass[10pt,letterpaper]{article}
\usepackage{technical-user-manual}
\setDocTitle{Computer Security Engineering Curriculum}
\setDocSubtitle{36-week prerequisite-based pathway}
\setDocVersion{1.0}
\setDocStatus{Independent educational companion}
\setDocOwner{latex-docs / security}
\setDocDate{2026-11-02}
\setDocSummary{Primary 2019 textbook anchors and independent exercises with gated evidence.}
\begin{document}
\makedoctocpage
\begin{notebox}[Source and authority boundary]
\textbf{BOOK:} Matt Bishop, \emph{Computer Security: Art and Science}, Second Edition (Pearson, 2019). Book-aligned sections are \emph{original summaries and reading prompts}, not copied textbook content.
\textbf{EXTENSION:} Any contemporary framework mappings, implementation tasks, scoring rules, and operational recommendations are independently authored educational supplements, not the book's positions. Nothing here proves compliance.
\end{notebox}
\section{Entry requirements and diagnostic}
Required: reading technical English, Git, basic Linux shell, programming in Python and C or C++, elementary sets/graphs, and the ability to run tests without production privileges. Diagnostics: explain a trust boundary, write a unit test for rejected input, model a small directed graph, and distinguish a proof from a test.
\section{Prerequisite dependency graph}
\textbf{Order:} CIA and policy $\rightarrow$ protection-state models $\rightarrow$ confidentiality/integrity $\rightarrow$ information flow $\rightarrow$ identity and access control $\rightarrow$ assurance and evaluation. Cryptography needs key lifecycle and threat modeling. Operations require evidence provenance and authority to test. Learners may study nondependent units in parallel only after their prerequisites.
\section{36-week curriculum and evidence gates}
\begin{longtable}{@{}p{0.10\linewidth}p{0.19\linewidth}p{0.12\linewidth}p{0.23\linewidth}p{0.25\linewidth}@{}}
\toprule Weeks & Module & Chapters & Prerequisites and skills & Gate deliverable\\ \midrule\endhead
1-2 & Foundations & 1, 4 & CIA, policy precision, adversary assumptions & Threat surface register and policy decision table\\ \midrule
3-5 & Protection state and decidability & 2, 3 & Sets, graphs, state transitions, finite search & Access matrix simulator and safety limitation statement\\ \midrule
6-8 & Confidentiality and integrity & 5, 6 & Partial orders and lattice concepts & Bell-LaPadula/Biba/Clark-Wilson rule comparisons\\ \midrule
9-10 & Availability and history policies & 7, 8 & Queues, state machines and resource accounting & Bounded load model and conflict-of-interest test\\ \midrule
11-13 & Noninterference and channels & 9, 17, 18 & Deterministic programs, observational equivalence & Two-run security tests and channel inventory\\ \midrule
14-17 & Cryptography and identity & 10-13, 15 & Discrete mathematics and API threat boundaries & Authenticated encryption and credential lifecycle plan\\ \midrule
18-20 & Systems enforcement & 14, 16, 28, 29 & Operating systems and network architecture & Central enforcement and least-privilege architecture review\\ \midrule
21-24 & Assurance and evaluation & 19-22 & Logic, test design, proof obligations & Assurance case and bounded model checking report\\ \midrule
25-28 & Threats, observability and response & 23-27 & Logs, incident handling, vulnerability semantics & Synthetic detection and triage case study\\ \midrule
29-31 & Human and program security & 30, 31 & C/C++ and Python basic secure coding & Secure parser test evidence and user recovery review\\ \midrule
32-36 & Integrated capstone & All relevant chapters & Prior gates passed & End-to-end assurance package with independent peer review\\ \midrule
\bottomrule
\end{longtable}
\section{Progression rules}
\begin{enumerate}
\item Complete a pretest covering model terminology and safety restrictions. Incorrect or unavailable prerequisites trigger remedial readings and repeat exercises.
\item Each module needs a reviewable design, executable or logically checkable exercise, one negative test, and a limitations statement.
\item Peer review must distinguish intended policy, actual mechanism, evidence and residual risk. Record a decision before advancing.
\item Modules 3, 5 and 8 are formal security gates: learners must articulate assumptions, invariant and counterexample conditions.
\item Capstone approval is educational only and is not an audit opinion, certification or authorization to operate.
\end{enumerate}
\section{Suggested capstone}
Model a multi-tenant document API with a policy matrix, identity binding, confidentiality constraints, integrity workflow, immutable event schema, bounded resource access, and recovery assumptions. Implement isolated Python/C/C++ components, run negative tests, identify one counterexample, build an assurance case and record risk ownership. Use synthetic data and non-production services exclusively.
\section{Rubric}
Score each domain 0--3: security objectives and scope; formal precision; executable evidence; negative-path coverage; provenance/reproducibility; residual risk; and peer review. A total score is diagnostic only; any unreviewed critical policy bypass is a separate gate failure.

\end{document}
